% !TEX root = ../journal_0421052099.tex
% ---------------------------------------------------------------------
% Reframed in September 2026: the attribution study is the contribution
% and the system is the instrument. An earlier draft led with the system
% and its verification layer, and a reviewer pass found that framing
% could not be defended -- the layer changes the answer on one question
% per dataset, zero ungrounded assertions belong to navigation, and the
% output contract is not persisted.
%
% What this section must still not do: promise that verification raises
% accuracy, or say "auditable" without saying where. The layer is
% described at its real width -- relation-blind structural routes -- and
% Sections 3 and 7 bound it.
%
% Register: short sentences, one idea each, plain connectives, no noun
% stacks. Filled at 72 columns.
%
% CUT 2026-09-23 to ~430 words, from ~960. Gone: the five-item
% \enumerate of contributions and the section-by-section roadmap, which
% between them said three times what the findings paragraph says once.
% Every number and every \Cref that check_paper_numbers.py pins
% survived.
% ---------------------------------------------------------------------

\section{Introduction}\label{sec:intro}

Large language models answer factual questions fluently, and often
enough they answer them wrongly. A fabricated answer reads in the same
register as a correct one, so it is hard to catch from the output
alone~\cite{huang2025survey}, and the problem compounds once an answer
requires several facts put together: if the first hop lands on the wrong
entity, nothing later in the chain will know. Retrieval-augmented
generation~\cite{lewis2020retrieval} grounds the model in retrieved
text, but it retrieves once, before any reasoning starts, and the
passage that holds an intermediate entity need not resemble the question
at all. A knowledge graph~\cite{bollacker2008freebase} turns composition
into traversal, since each hop is an edge, and static graph
retrieval~\cite{edge2024local} still fetches a neighbourhood in one
shot. Agentic systems close that gap by interleaving retrieval with
reasoning: the model takes one step, looks at what it found, and decides
where to look
next~\cite{sun2024think,chen2024plan,luo2024reasoning,luo2025kbqao1}.

These systems are assemblies. A typical one decomposes the question into
sub-goals, scores candidate edges with the model, backtracks from dead
ends, and checks its own output, and the literature reports the assembly
as a whole. When it outperforms a static baseline, three things could
have produced the gain: a component, the budget the comparison ran
under, or the labels of the benchmark itself. The report rarely
separates them, so later work inherits the whole assembly on faith,
along with its cost. The exception shows why taking a system apart
matters. KBQA-o1's own ablation removes its tree search, and GrailQA F1
falls by thirty points~\cite{luo2025kbqao1}. This paper does the same
for an agentic knowledge-graph question-answering system built so that
each of its components can be removed cleanly, and it asks which of the
three sources produced the gain.

The system is \textbf{Agentic Graph Reasoning} (AGR), an explicit state
machine over a constrained, deterministic graph tool API: a program runs
the loop, and the language model is consulted only at bounded decision
points. Four components can be removed. A \emph{planner} decomposes the
question into an ordered chain of sub-objectives. A hybrid \emph{edge
scorer} blends embedding similarity with a batched model judgement. A
\emph{backtracker} restores an earlier frontier and bans the edges that
led away from it. And a \emph{structural verification layer} splits each
draft answer into atomic claims and checks them against the traversed
triples before the answer is emitted. The last is the component the
system was designed around, so we should be clear about what it actually
checks: its two graph routes ask whether an edge joins a claim's
endpoints, not whether it is the edge the claim asserts, so a claim that
X is Y's mother is certified by an edge recording that X is Y's child.
\Cref{sec:verification,sec:limits-contribution} bound it further.

On $400$ questions from each of WebQSP~\cite{yih2016value} and
ComplexWebQuestions~\cite{talmor2018web}, over a Freebase-derived
environment of $2.59$ million entities and $8.31$ million triples, AGR
reaches Hits@1 of $0.755$ and $0.522$ and leads four baselines that
share its backbone and its call budget. That is where the analysis
starts, not where it ends. Against a reimplementation of
Think-on-Graph~\cite{sun2024think} on the same tools, the entire margin
comes from the $29\%$ and $44\%$ of questions on which the shared
$25$-call cap truncates the baseline; where the cap lets it finish, it
is ahead (\Cref{sec:margin}). The ablation finds one component with a
detectable effect, and the effect runs the wrong way. Removing the
planner improves WebQSP F1 by $0.083$ while cutting tokens by $31\%$,
because decomposing a question that has no compositional structure
throws away the context that would have answered it. Backtracking, model
scoring, and claim verification show no effect at the power available,
and \Cref{sec:power} states the smallest effect each test could have
caught. The failure census finds label defects on $5.5\%$ and $4.8\%$ of
the sampled questions, and names a failure that passes every grounding
check and that consensus-based label repair cannot see.
